\chapter*{ЗАКЛЮЧЕНИЕ}
\addcontentsline{toc}{chapter}{ЗАКЛЮЧЕНИЕ}

Разработана программа для предварительного отбора JPEG-файлов с признаками повторного сжатия. Она принимает фотографию или каталог и выдаёт класс и оценку двукратного сжатия. Программа извлекает яркостной фрагмент и таблицу квантования непосредственно из конечного JPEG.

По результатам анализа опубликованных методов выбрано общее представление изображения для двух классификаторов. На 30\,000~размеченных JPEG обучены многослойный персептрон и одномерная свёрточная сеть. Для применения выбрана свёрточная сеть. На тесте из 5\,000~файлов доля верных ответов составила 80,44\,\% против 78,72\,\% у персептрона. Сеть дала на 92~ложных срабатывания меньше и пропустила на шесть повторно сжатых файлов больше. При загруженной модели медианное время обработки фрагмента $256\times256$ на центральном процессоре составило 0,546~мс.

Эксперимент на 300~отдельных фотографиях UCID определил условия, при которых такой отбор полезен. При совпадающей сетке блоков и повышении качества второго сохранения программа обнаружила 75,33\,\% повторно сжатых файлов. При одинаковых значениях качества полнота составила 1,27\,\%, а после сдвига сетки на четыре пикселя --- 1,21\,\% по всем парам настроек. Поэтому файлы с положительным ответом можно передать на дальнейший анализ, а отрицательный ответ при неизвестной истории обработки оставляет вопрос о повторном сохранении открытым.
